\textbf{COLUNI 2026 -- Questão 09}\par\medskip Leia uma parte do conto \textit{Com certeza tenho amor}, de Marina Colassanti, e observe a estrutura sintática:

\begin{quote}“Moça tão resguardada por seus pais não deveria ter ido à feira. Nem foi, embora muito o desejasse. Mas porque o desejava, convenceu a ama que a acompanhava a tomar uma rua em vez de outra para ir à igreja, e a rua que tomaram passava tão perto da feira que seus sons a percorriam como água e as cores todas da feira pareciam espelhar-se nas paredes claras. Foi dessa rua, olhando através do véu que lhe cobria metade do rosto, que a moça viu os saltimbancos em suas acrobacias. E foi nessa rua, recortada como uma silhueta em suas roupas escuras, o rosto meio coberto por véu, que o mais jovem dos saltimbancos, atrasado a caminho da feira, a viu.” (p.41)\end{quote}

Na estruturação sintática da linguagem do excerto acima, há predominância de:
\begin{enumerate}
\item períodos simples e orações subordinadas, construindo um mundo simbólico de expectativas.
\item períodos simples e orações coordenadas, criando uma visão sobre a vida cotidiana.
\item períodos compostos com mais orações subordinadas, criando uma perspectiva do feminino.
\item períodos compostos com mais orações coordenadas, construindo uma concepção de vida urbana.
\end{enumerate}
